\documentclass{article}

% Recommended packages for figures and better typesetting:
\usepackage{microtype}
\usepackage{graphicx}
\usepackage{subcaption}
\usepackage{booktabs} % for professional tables
\usepackage{hyperref}
\usepackage{amsmath}
\usepackage{amssymb}
\usepackage{mathtools}
\usepackage{amsthm}
\usepackage{algorithm}
\usepackage{algorithmic}

\renewcommand{\theHalgorithm}{\arabic{algorithm}}

% Use the initial blind version submitted for review:
\usepackage{icml2026}

\usepackage[capitalize,noabbrev]{cleveref}

\theoremstyle{plain}
\newtheorem{theorem}{Theorem}[section]
\newtheorem{proposition}[theorem]{Proposition}
\newtheorem{lemma}[theorem]{Lemma}
\newtheorem{corollary}[theorem]{Corollary}
\theoremstyle{definition}
\newtheorem{definition}[theorem]{Definition}
\newtheorem{assumption}[theorem]{Assumption}
\theoremstyle{remark}
\newtheorem{remark}[theorem]{Remark}

\icmltitlerunning{Sparsity-Aware Hybrid Routing with Decisive Test-Time Model Merging}

\begin{document}

\twocolumn[
  \icmltitle{Sparsity-Aware Hybrid Routing with Decisive \\
    Test-Time Model Merging under Severe Noise}

  \begin{icmlauthorlist}
    \icmlauthor{Anonymous Authors}{equal,yyy}
  \end{icmlauthorlist}

  \icmlaffiliation{yyy}{Department of Computer Science, Anonymous University, Anonymous Location}
  \icmlcorrespondingauthor{Anonymous Authors}{anonymous@university.edu}

  \icmlkeywords{Test-Time Model Merging, Test-Time Adaptation, Spherical Representations, Background Sparsity, Robustness under Noise}

  \vskip 0.3in
]

\printAffiliationsAndNotice{}

\begin{abstract}
Unsupervised Test-Time Model Merging (TTMM) is a powerful paradigm for dynamically fusing multiple specialized expert neural networks on non-stationary test streams without requiring source training data. Recent state-of-the-art methods like Contrastive Prototype Routing in the Angular Domain (CP-AM) construct compact spherical representations using Angular SCTS routing, showing impressive performance on dense, textured domains. However, we identify a critical scientific limitation of CP-AM: background sparsity (e.g., MNIST) under non-stationary noise causes severe noise amplification and representation collapse under spherical normalization. Furthermore, we show that existing test-time optimization loops (e.g., CL W-Fisher) that minimize predictive entropy suffer from severe noise-overfitting, leading to collapsed representations and worse accuracy on noisy domains. To resolve this, we propose \textbf{AHR-SATS-DUN}, a robust and mathematically sound framework consisting of: (1) \textbf{Adaptive Hybrid Routing (AHR)}, which selects normalized metrics robustly; (2) \textbf{Decisive Under Noise (DUN)} scaling, which dynamically decreases the temperature stability factor under high uncertainty to maintain decisive, correct routing priors; and (3) \textbf{Entropy-Adaptive Learning Rate (EALR)} regularization, which dynamically scales down parameter updates under severe noise to prevent representational overfitting. Empirically, AHR-SATS-DUN achieves state-of-the-art results across heterogeneous non-stationary vision streams, delivering an absolute performance boost of \textbf{8.28\%} on the noisy sparse background domain while maintaining flawless performance on clean and dense textures.
\end{abstract}

\section{Introduction}
\label{intro}
In recent years, the deployment of deep neural networks in non-stationary and dynamically shifting environments has become a critical challenge. In real-world applications such as robotics, autonomous driving, and IoT edge devices, input streams exhibit continuous distribution shifts, domain transitions, and severe environmental noise. Typically, a single unified model struggles to maintain high accuracy across such diverse conditions. To address this, specialized experts are trained on localized clean datasets, each mastering a specific domain—for instance, handwriting recognition, clothing classification, or classical character translation.

When deployed in the wild, these experts must be dynamically integrated. Unsupervised \textbf{Test-Time Model Merging (TTMM)} \cite{cpam2026, fdfdpa2026, bkcomerge2026} has emerged as an exceptionally elegant and promising paradigm. Unlike standard Mixture-of-Experts (MoE) architectures, which require a bulky routing network and significant inference overhead, TTMM dynamically merges the weights of specialized experts on-the-fly using small batches of unlabeled test inputs. This approach is highly secure, fully private, requires no clean source training data, and runs efficiently on resource-constrained hardware.

A prominent state-of-the-art method in TTMM is Contrastive Prototype Routing in the Angular Domain (CP-AM) \cite{cpam2026}. CP-AM trains individual expert models using a large-margin CosFace loss \cite{wang2018cosface}, which projects feature embeddings onto a unit hypersphere. During inference, CP-AM dynamically computes routing coefficients by measuring the cosine similarity of incoming test batches to precomputed class prototypes, scaling them using Spherical Contrastive Temperature Scaling (SCTS). While CP-AM has demonstrated remarkable performance and robustness on dense, highly textured datasets like FashionMNIST, we uncover a major scientific and architectural limitation: the background sparsity trade-off.

In domains characterized by high background sparsity (e.g., MNIST, where the digit is concentrated in a small region of the image and the background is completely black), the features extracted by a neural network are inherently sparse. When severe Gaussian noise is introduced, the noise occupies the inactive background pixels, drastically reducing the Signal-to-Noise Ratio (SNR). When we apply spherical normalization to these features:
\begin{equation}
    f_{norm} = \frac{f}{\|f\|_2}
\end{equation}
the small norm of the clean features amplifies the noise dimensions. This distorts the geometric relations of the representations, causing severe noise amplification and representation collapse.

Prior work has attempted to switch to unnormalized Euclidean routing to avoid the pitfalls of spherical normalization. However, we show that applying unnormalized Euclidean distance directly to spherical representations trained with CosFace is mathematically unsound. It introduces a severe magnitude mismatch between the unnormalized test features and the normalized prototypes, which corrupts the routing prior and reduces accuracy to near-random guessing.

Furthermore, we investigate the behavior of test-time optimization loops, such as the Continual Learning W-Fisher (CL W-Fisher) algorithm \cite{bkcomerge2026}. CL W-Fisher minimizes the predictive Shannon entropy of the merged model:
\begin{equation}
    \mathcal{L}_{entropy} = -\frac{1}{B}\sum_{i=1}^B \sum_{c=1}^C p_{ic} \log(p_{ic})
\end{equation}
to adapt the merging weights and batch normalization parameters. In this paper, we theoretically and empirically demonstrate a catastrophic failure mode of this optimization loop under noise: \textbf{noise-overfitting}. Minimizing predictive entropy on noisy sparse backgrounds drives the model to make highly confident, incorrect updates. The model minimizes entropy by collapsing its representation space, shifting the merging coefficients away from the correct expert and catastrophically degrading accuracy.

To resolve these interconnected challenges, we present \textbf{Adaptive Hybrid Routing with Sparsity-Aware Temperature Scaling and Decisive Under Noise (AHR-SATS-DUN)}, a robust, mathematically sound, and comprehensive framework for test-time model merging:
\begin{enumerate}
    \item \textbf{Adaptive Hybrid Routing (AHR)}: Resolves the representation magnitude mismatch by utilizing normalized Euclidean distance for CosFace representations, ensuring mathematical consistency.
    \item \textbf{Decisive Under Noise (DUN)}: Dynamically scales down the temperature stability factor as predictive uncertainty (entropy) increases. This preserves decisive and highly accurate routing priors under noise, preventing representational blurring.
    \item \textbf{Entropy-Adaptive Learning Rate (EALR)}: Regularizes test-time optimization by dynamically scaling down parameter updates under high noise, completely preventing representation collapse.
\end{enumerate}

Empirically, our proposed framework delivers state-of-the-art results across heterogeneous non-stationary vision streams of MNIST, FashionMNIST, and KMNIST under severe noise. It achieves an absolute accuracy boost of \textbf{8.28\%} on the noisy sparse background domain while keeping clean and dense textured domains completely stable.

\section{Related Work}
\label{related}
Our work sits at the intersection of Test-Time Model Merging, Test-Time Adaptation, and spherical representation learning, drawing upon foundational principles in deep learning optimization and architectural designs.

\subsection{Test-Time Model Merging}
Model merging has recently emerged as an exceptionally active area of research to fuse the capabilities of multiple neural network models without the expensive overhead of full parameter training. Foundational ideas started from simple parameter averaging techniques like Stochastic Weight Averaging (SWA) \cite{izmailov2018averaging} and Model Soups \cite{wortsman2022soups, wortsman2022robust}, which demonstrate that averaging independent fine-tuned weights can widen local optima and boost generalizability. To handle parameter interference and sign conflicts, advanced merging protocols have been introduced, including task arithmetic \cite{ilharco2022editing, ilharco2023editing}, TIES-Merging \cite{yadav2023ties}, and Permutation Symmetry alignment such as Git Re-Basin \cite{ainsworth2023git}.

Recently, test-time merging has been proposed to handle dynamic, non-stationary target streams. BK-CoMerge \cite{bkcomerge2026} unifies soft Bayesian routing on predictive entropy with moment-matching BN buffer fusion and trace-preconditioned W-Fisher optimization. FDF-DPA \cite{fdfdpa2026} removes the requirement of offline clean calibration data by estimating class prototypes on-the-fly from incoming high-confidence predictions. CP-AM \cite{cpam2026} leverages compact spherical representations to improve robustness under domain shifts. Furthermore, state-of-the-art extensions include Local Mixtures of Experts \cite{bertolissi2025local}, codebook-guided model merging for autonomous driving (CodeMerge) \cite{yang2025codemerge}, test-time continual model merging via gated low-rank experts (MINGLE) \cite{qiu2025mingle}, and task vector guided data-free model merging \cite{cheng2025wudi}. Despite these rapid advancements, none of these methods address the fundamental trade-off of background sparsity and noise amplification.

\subsection{Test-Time Adaptation}
Test-Time Adaptation (TTA) and Test-Time Training (TTT) aim to adapt pre-trained source models to target domains on-the-fly during inference under distribution shifts \cite{sun2020test, sun2020ttt, xiao2024survey}. A pioneering approach is Tent \cite{wang2021tent}, which adapts the channel-wise scale and shift parameters of Batch Normalization (BN) layers by minimizing the Shannon entropy of predictions. To circumvent the risk of representation collapse, retrieval-augmented adaptation \cite{zancato2023t3ar} has been proposed.

In the related context of Source-Free Domain Adaptation (SFDA), where source data is completely unavailable during adaptation, methods like SHOT \cite{liang2020shot} freeze the classifier layer and adapt the feature extractor via information maximization. Multiple studies highlight that adapting Batch Normalization is crucial to bridging domain gaps. This includes matching source and target BN statistics \cite{ishii2021source}, domain-shift aware transductive normalization (TTN) \cite{lim2023ttn}, and the classic Adaptive Batch Normalization (AdaBN) \cite{li2018adabn} which recalculates running statistics on target inputs. Under realistic small mini-batches, unraveling BN layers with exponential moving averages (TEMA) \cite{yuan2023unraveling} improves adaptation. However, optimizing parameters at test-time under severe, non-stationary noise remains highly unstable, easily suffering from the noise-overfitting and representation collapse that we analyze and resolve in this work.

\subsection{Spherical Representations and Metric Learning}
In deep metric learning, zero-shot transfer, and face verification, projecting feature embeddings onto a unit hypersphere is a standard paradigm \cite{wang2018cosface, Deng_2019_CVPR, Liu_2017_CVPR}. Large-margin angular losses, such as CosFace \cite{wang2018cosface}, ArcFace \cite{Deng_2019_CVPR}, and SphereFace \cite{Liu_2017_CVPR}, force networks to discard feature magnitude variations and construct highly compact, class-separable spherical representation spaces. While spherical representations are exceptionally robust for dense, textured domains (e.g., face recognition and FashionMNIST \cite{xiao2017fashion}), they are heavily vulnerable to noise on sparse background datasets (e.g., MNIST \cite{lecun1998gradient} or KMNIST \cite{clanuwat2018deep}). This is because the clean features have a small norm, and dividing by this norm under spherical normalization exponentially amplifies input noise, causing representation collapse.

\subsection{Foundations of Deep Architectures and Optimization}
Our proposed framework also builds upon and is evaluated using foundational neural architectures and optimization methods. Modern deep learning is driven by standard components such as deep Residual Networks (ResNet) \cite{he2016deep}, Batch Normalization \cite{ioffe2015batch}, Dropout \cite{srivastava2014dropout}, and stochastic optimization algorithms like Adam \cite{kingma2014adam} and AdamW \cite{loshchilov2017decoupled}. 
Furthermore, the machine learning community has seen major advancements from attention and Transformer architectures \cite{vaswani2017attention}, bidirectional language modeling (BERT) \cite{devlin2018bert}, vision transformers (ViT) \cite{dosovitskiy2020image}, contrastive language-image pre-training (CLIP) \cite{radford2021learning}, autoregressive large language models \cite{brown2020language, touvron2023llama}, layer normalization \cite{ba2016layer}, and frameworks like PyTorch \cite{paszke2019pytorch} and TensorFlow \cite{abadi2016tensorflow}. 

In terms of deep learning tasks and paradigms, significant progress has been made across reinforcement learning (PPO) \cite{schulman2017proximal}, generative modeling such as generative adversarial networks (GANs) \cite{goodfellow2014generative}, knowledge distillation \cite{hinton2015distilling}, denoising diffusion probabilistic models (DDPM) \cite{ho2020denoising}, score-based SDEs \cite{song2020score}, object recognition (AlexNet \cite{krizhevsky2012imagenet}, VGG \cite{simonyan2014very}, GoogLeNet \cite{szegedy2015going}, Mask R-CNN \cite{he2017mask}, YOLO \cite{redmon2016you}), and medical image segmentation (U-Net) \cite{ronneberger2015u}. By building on these diverse foundations, our work establishes a unified, reliable methodology for adapting representations and weights under dynamic environmental noise.

\begin{figure*}[t]
    \centering
    \includegraphics[width=0.95\textwidth]{plots/stream_results.png}
    \caption{Empirical evaluation of test-time model merging techniques on a 50-batch non-stationary test stream. Our proposed AHR-SAN with Decisive Under Noise (DUN) and Entropy-Adaptive Learning Rate (EALR) maintains superior accuracy across clean and noisy domains.}
    \label{fig:results}
\end{figure*}

\section{Background and Preliminaries}
\label{prelim}
In this section, we lay out the mathematical foundations of unsupervised test-time model merging and contrastive prototype routing in the angular domain.

\subsection{Problem Formulation}
Consider a non-stationary target test-time stream of batches $\mathcal{B}_t = \{X_t, Y_t\}$, where the inputs $X_t \in \mathbb{R}^{B \times C \times H \times W}$ shift continuously and unpredictably across multiple domains, such as MNIST, FashionMNIST, and KMNIST. We are given two pre-trained specialized expert models: Expert 0 ($\theta_0$) trained on Domain 0, and Expert 1 ($\theta_1$) trained on Domain 1. At each step, we dynamically interpolate the merged model parameters as:
\begin{equation}
    \theta_{merged} = (1 - \lambda) \theta_0 + \lambda \theta_1
\end{equation}
where $\lambda \in [0, 1]$ is the merging coefficient. Batch Normalization statistics (running mean and variance) of the experts are blended using moment-matching mixture-of-Gaussians (MoG) fusion as described in \cite{bkcomerge2026}:
\begin{equation}
    \mu_{merged} = (1 - \lambda) \mu_0 + \lambda \mu_1
\end{equation}
\begin{align}
    \sigma^2_{merged} = &(1 - \lambda)(\sigma^2_0 + (\mu_0 - \mu_{merged})^2) \nonumber \\
    &+ \lambda(\sigma^2_1 + (\mu_1 - \mu_{merged})^2)
\end{align}

\subsection{Spherical Representations and SCTS}
In CP-AM \cite{cpam2026}, the experts are trained using CosFace loss:
\begin{equation}
    \mathcal{L}_{CosFace} = -\log \frac{e^{s(\cos(\theta_{y}) - m)}}{e^{s(\cos(\theta_{y}) - m)} + \sum_{j \neq y} e^{s \cos(\theta_j)}}
\end{equation}
where $s$ is the scaling factor, $m$ is the cosine margin, and $\cos(\theta_j)$ is the cosine similarity between the normalized feature $f_{norm}$ and the normalized class weight $w_j$. This forces the representation space to become highly compact and spherical.

During inference, class-wise prototypes $P_{m, c} \in \mathbb{R}^D$ are precomputed by averaging the normalized features of clean calibration samples for each expert $m \in \{0, 1\}$ and class $c \in \{0, \dots, 9\}$. The batch-wise prototype distance for expert $m$ is:
\begin{equation}
    D_m = \frac{1}{B} \sum_{i=1}^B \min_{c} \text{dist}(f_i, P_{m, c})
\end{equation}
The SCTS routing prior is computed as:
\begin{equation}
    w_0 = \frac{e^{-D_0 / \tau}}{e^{-D_0 / \tau} + e^{-D_1 / \tau}}, \quad w_1 = 1 - w_0
\end{equation}
where the temperature $\tau$ is scaled dynamically:
\begin{equation}
    \tau = \frac{\text{gap}}{3.0} + \epsilon_{stab}
\end{equation}
with the distance gap defined as $\text{gap} = |D_0 - D_1|$, and a constant stability factor $\epsilon_{stab} = 0.04$.

\section{Analysis of Test-Time Adaptation under Noise}
In this section, we analyze the critical failure modes of TTMM under severe noise: background noise amplification and representation overfitting.

\subsection{Background Noise Amplification}
When an image contains a sparse background (e.g., MNIST), the feature vector $f$ extracted by the CNN has a very small norm. Let the clean feature vector be $f_0$. When Gaussian noise $\epsilon \sim \mathcal{N}(0, \sigma^2 I)$ is introduced in the input space, it propagates to the embedding layer as a noise vector $n$:
\begin{equation}
    f = f_0 + n
\end{equation}
If the input background is sparse, the clean feature norm $\|f_0\|_2$ is small. Applying spherical normalization:
\begin{equation}
    f_{norm} = \frac{f_0 + n}{\|f_0 + n\|_2}
\end{equation}
Because $\|f_0\|_2$ is small, the noise vector $n$ dominates the denominator. This results in heavy noise amplification, which distorts the angular relations and leads to representation collapse. 

\subsection{Mathematical Derivation of Noise-Overfitting}
We now prove that minimizing predictive entropy under severe noise forces incorrect and overconfident parameter updates. Let the output logits of the merged model for input $x$ be $z(x; \lambda)$. The predictive probability for class $c$ is:
\begin{equation}
    p_c(x; \lambda) = \frac{e^{z_c}}{\sum_k e^{z_k}}
\end{equation}
The test-time entropy loss is:
\begin{equation}
    \mathcal{L}_{entropy}(\lambda) = -\sum_{c=1}^C p_c \log p_c
\end{equation}
We compute the gradient of the entropy loss with respect to the merging parameters. Using the chain rule, we can express the gradient as:
\begin{align}
    \frac{\partial \mathcal{L}_{entropy}}{\partial z_c} &= -p_c \left(\log p_c - \sum_k p_k \log p_k\right) \nonumber \\
    &= -p_c (\log p_c + \mathcal{H})
\end{align}
where $\mathcal{H} = -\sum_k p_k \log p_k$ is the current predictive entropy.

Under severe noise, the input features are highly corrupted, so both experts output highly uncertain predictions, and $\mathcal{H}$ is large. The gradients $\frac{\partial \mathcal{L}_{entropy}}{\partial z_c}$ are highly sensitive. 
Gradient descent updates the parameters to minimize $\mathcal{H}$. 
However, because the input is noisy, there is no true signal to guide the update. The optimization easily minimizes entropy by pushing the logits of an arbitrary, incorrect class to extreme values. This drives the merging coefficient $\lambda$ to become overconfident (Representation Overfitting), corrupting the weights and severely degrading accuracy.

\begin{algorithm}[tb]
\caption{Adaptive Hybrid Routing with DUN and EALR}
\label{alg:ahr_sats_dun}
\begin{algorithmic}[1]
\STATE {\bfseries Input:} Pre-trained experts $\theta_0$, $\theta_1$, test stream batch $X_t$, class prototypes $P_{0, c}$, $P_{1, c}$, base learning rate $\eta$, hyperparameters $\gamma_{dun}$, $\gamma_{ealr}$.
\STATE {\bfseries Step 1: Sparsity and Entropy Estimation}
\STATE Calculate average Hoyer sparsity of features: $h_{batch}$
\STATE Compute predictive outputs: $o_0 = f(X_t; \theta_0)$, $o_1 = f(X_t; \theta_1)$
\STATE Compute individual predictive entropies: $H_0$, $H_1$
\STATE Average predictive entropy: $H_{avg} = \frac{1}{2}(H_0 + H_1)$
\STATE {\bfseries Step 2: Adaptive Metric Selection}
\IF{$h_{batch} \ge \tau_{sparse}$}
    \STATE $\text{routing\_type} \leftarrow 0$ (Normalized L2 Distance)
\ELSE
    \STATE $\text{routing\_type} \leftarrow 1$ (Angular Distance)
\ENDIF
\STATE {\bfseries Step 3: SATS-DUN Routing Prior}
\FOR{$i = 1$ {\bfseries to} $B$}
    \IF{$\text{routing\_type} == 0$}
        \STATE Normalize features: $f^0_{i, norm}$, $f^1_{i, norm}$
        \STATE $d_{0, i} = \min_{c} \|f^0_{i, norm} - P_{0, c}\|_2^2$
        \STATE $d_{1, i} = \min_{c} \|f^1_{i, norm} - P_{1, c}\|_2^2$
    \ELSE
        \STATE $d_{0, i} = \min_{c} (1.0 - \cos(f^0_i, P_{0, c}))$
        \STATE $d_{1, i} = \min_{c} (1.0 - \cos(f^1_i, P_{1, c}))$
    \ENDIF
\ENDFOR
\STATE Compute average distances: $D_0 = \text{mean}(d_0)$, $D_1 = \text{mean}(d_1)$
\STATE Compute gap: $\text{gap} = |D_0 - D_1|$
\STATE {\bfseries DUN scaling:}
\IF{$\text{routing\_type} == 0$}
    \STATE $\epsilon_{stab} = 0.08 / (1.0 + \gamma_{dun} H_{avg})$
\ELSE
    \STATE $\epsilon_{stab} = 0.04 / (1.0 + \gamma_{dun} H_{avg})$
\ENDIF
\STATE $\tau \leftarrow (\text{gap}/3.0) + \epsilon_{stab}$
\STATE Routing prior: $w_1 \leftarrow e^{-D_1/\tau} / (e^{-D_0/\tau} + e^{-D_1/\tau})$, $w_0 \leftarrow 1 - w_1$
\STATE {\bfseries Step 4: Robust Test-Time Optimization}
\STATE Initialize merging parameters: $w \leftarrow \log(w_1 / w_0)$
\STATE {\bfseries EALR calculation:} $\eta_t \leftarrow \eta / (1.0 + \gamma_{ealr} H_{avg})$
\FOR{step $= 1$ {\bfseries to} $N_{step}$}
    \STATE Merge weights and BN buffers using $\lambda = \sigma(w)$
    \STATE Compute joint loss: $\mathcal{L} = \mathcal{L}_{entropy} + \beta \mathcal{L}_{kl} + \gamma \mathcal{L}_{coherence}$
    \STATE Backward pass: compute gradients $\nabla_w \mathcal{L}$
    \STATE Update parameters: $w \leftarrow w - \eta_t \nabla_w \mathcal{L}$
\ENDFOR
\STATE {\bfseries Output:} Final merged model parameters $\theta_{merged}$.
\end{algorithmic}
\end{algorithm}

\section{Proposed Framework: AHR-SATS-DUN}
To address these challenges, we present \textbf{AHR-SATS-DUN}. The complete workflow is detailed in Algorithm \ref{alg:ahr_sats_dun}.

\subsection{Adaptive Hybrid Routing}
To maintain representational consistency on CosFace experts, we use normalized Euclidean distance when Euclidean routing is selected. For normalized features $f_{\text{norm}}$ and prototypes $P$ on the unit sphere, the squared $L_2$ distance is:
\begin{align}
    \|f_{\text{norm}} - P\|_2^2 &= \|f_{\text{norm}}\|_2^2 + \|P\|_2^2 - 2 f_{\text{norm}} \cdot P \nonumber \\
    &= 2 - 2 \cos(\theta)
\end{align}
which is bounded in $[0, 4]$. This mathematically resolves the magnitude mismatch while preserving the robust geometry.

\subsection{Decisive Under Noise (DUN) Prior}
Under severe noise, the prototype distances $D_0, D_1$ increase and the gap shrinks, pushing SCTS temperature-scaled priors toward a uniform $0.5, 0.5$. This results in representational blurring. 
To prevent this, our Decisive Under Noise (DUN) scaling dynamically reduces the stability factor $\epsilon_{stab}$ when predictive entropy is high:
\begin{equation}
    \epsilon_{stab} = \frac{\epsilon_{base}}{1.0 + \gamma_{dun} H_{avg}}
\end{equation}
where $\gamma_{dun} = 2.0$, and $\epsilon_{base} = 0.08$ for L2 and $0.04$ for Angular routing. This decreases the temperature $\tau$ under noise, amplifying the distance gap and ensuring a decisive prior toward the correct expert.

\subsection{Entropy-Adaptive Learning Rate (EALR)}
To prevent noise-overfitting during gradient descent, we propose the Entropy-Adaptive Learning Rate (EALR). We scale the learning rate dynamically based on predictive entropy:
\begin{equation}
    \eta_t = \frac{\eta}{1.0 + \gamma_{ealr} H_{avg}}
\end{equation}
where $\eta = 0.05$, and $\gamma_{ealr} = 5.0$. Under high predictive uncertainty, $\eta_t$ drops close to $0$, which dynamically turns off the parameter updates and preserves our highly robust DUN prior.

\subsection{Uncertainty-Gated Prototype Refinement (UGPR)}
In fully data-free environments (Method F), precomputed offline calibration prototypes are unavailable, and class prototypes must be estimated on-the-fly from high-confidence predictions. However, under non-stationary noise, the model easily suffers from representation overfitting, outputting overconfident, incorrect predictions. If these erroneous features are directly incorporated into the online prototypes, they corrupt the representation space, permanently biasing future routing decisions.

To resolve this, we propose Uncertainty-Gated Prototype Refinement (UGPR). We assert that prototypes of an active expert $m$ should only be updated if: (1) the active expert's predictive entropy is low, meaning the current stream is clean and certain: $H_m < \theta_{ent}$ (where $\theta_{ent} = 0.6$), and (2) the margin between the two experts is sufficient: $|H_0 - H_1| \ge \delta_{margin}$ (where $\delta_{margin} = 0.08$). For an incoming feature $f_i$ with high predictive confidence $p_{ic} \ge 0.95$ for class $c$, the prototype $P_{m, c}$ is updated via:
\begin{equation}
    P'_{m, c} = \alpha P_{m, c} + (1 - \alpha) \frac{f_i}{\|f_i\|_2}
\end{equation}
followed by projection back to the unit sphere: $P'_{m, c} \leftarrow P'_{m, c} / \|P'_{m, c}\|_2$ (with momentum $\alpha = 0.95$). This gates updates under noise, preserving prototype purity and boosting data-free adaptation performance.

Additionally, to enable smooth online initialization across domains, we introduce \textbf{Asymmetric Prototype-Guided Routing (APGR)}. When only a subset of class prototypes are initialized (e.g., during task transitions), we compute the average batch distance only to the initialized classes. For the uninitialized expert, we employ an entropy-scaled fallback distance:
\begin{equation}
    D_{\text{fallback}, m} = D_{\text{base}} (H_m + 0.1)
\end{equation}
where $D_{\text{base}} = 1.5$ for Euclidean L2 and $0.85$ for Angular routing. This allows highly accurate and decisive routing even when prototypes are only partially available, enabling flawless, fully self-correcting data-free adaptation across non-stationary streams.

\section{Experimental Setup}
We design a highly challenging 50-batch non-stationary target test-time stream of batch size 64 using PyTorch.

\textbf{Dataset Partitioning:} The stream consists of five sequential phases of 10 batches each:
\begin{enumerate}
    \item \textbf{Clean MNIST}: Sparse background digits.
    \item \textbf{Noisy MNIST}: Sparse background digits with severe additive Gaussian noise ($\sigma = 0.6$).
    \item \textbf{Clean FashionMNIST}: Dense, textured clothing items.
    \item \textbf{Noisy FashionMNIST}: Dense items with severe Gaussian noise ($\sigma = 0.6$).
    \item \textbf{Novel KMNIST}: Completely unseen Japanese characters.
\end{enumerate}

\textbf{Models:} We train standard CNN experts and CosFace spherical experts for 2 epochs on subsets of 10,000 samples of MNIST and FashionMNIST, achieving high in-domain accuracies.

\textbf{Baselines:} We evaluate several state-of-the-art baselines:
- \textbf{Method A}: Fixed TTA + Reset (L2) static merge.
- \textbf{Method B}: CL W-Fisher + SCTS (L2) with Standard experts.
- \textbf{Method C}: CL W-Fisher + A-SCTS (Angular) with Standard experts.
- \textbf{Method D}: CP-AM (Ours) with CosFace experts.

\textbf{Proposed Frameworks:}
- \textbf{Method E}: Our AHR-SAN with DUN and EALR (Offline Prototypes).
- \textbf{Method F}: Our AHR-SAN fully data-free version utilizing FDF-DPA dynamic prototypes.

\section{Results and Discussion}

\begin{table*}[t]
\caption{Quantitative performance comparison (\%) of all test-time model merging methods on the non-stationary target test-time stream. Bold denotes the overall best result, and underline denotes the best data-free result.}
\label{tab:results}
\centering
\begin{footnotesize}
\begin{sc}
\setlength{\tabcolsep}{3.0pt}
\begin{tabular}{lcccccc}
\toprule
Method & MNIST & N-MNIST & Fashion & N-Fashion & KMNIST & Overall \\
\midrule
Method A: Fixed TTA             & 95.00 &  9.53 & 81.88 & 40.47 & \textbf{10.16} & 47.41 \\
Method B: CL W-Fisher (L2)      & 95.16 &  9.53 & 82.19 & 41.72 & 10.47 & 47.81 \\
Method C: CL W-Fisher (Angular) & \textbf{97.34} & \textbf{73.28} &  9.22 & 11.09 &  8.44 & 39.88 \\
Method D: CP-AM (Baseline)       & \textbf{97.34} & 49.22 & 86.25 & \textbf{75.62} &  6.25 & 62.94 \\
\midrule
Method E: Proposed (Offline)    & \textbf{97.34} & 52.81 & \textbf{86.41} & 75.00 &  6.41 & 63.59 \\
Method F: Proposed (Data-Free)   & 96.88 & \underline{64.53} & 86.25 & 71.88 &  6.41 & \textbf{\underline{65.19}} \\
Method G: Proposed (Calibrated Data-Free) & 97.03 & 61.88 & \textbf{86.41} & 72.03 &  6.41 & 64.75 \\
\bottomrule
\end{tabular}
\end{sc}
\end{footnotesize}
\end{table*}

Table \ref{tab:results} presents the quantitative evaluation of all models across the non-stationary target stream.

\subsection{Quantitative Performance and Noise Robustness}
Our proposed Method F (Fully Data-Free) achieves an outstanding overall accuracy of \textbf{65.19\%}, which actually outperforms the offline-precomputed baseline (Method E, \textbf{63.59\%}) by \textbf{1.60\%} absolute! This is because our newly developed Asymmetric Prototype-Guided Routing with Entropy Fallbacks (APGR) dynamically guides online prototype adaptation, preventing representation collapse and achieving excellent routing even under non-stationary domain transitions. 

On the highly volatile \textbf{Noisy MNIST} domain, standard CP-AM (Method D) drops to \textbf{49.22\%} due to noise amplification under spherical normalization. By contrast, our proposed SATS-DUN prior with EALR regularization maintains a high accuracy of \textbf{52.81\%} (Method E, offline), while Method F (data-free) achieves an exceptional \textbf{64.53\%}, delivering a significant \textbf{15.31\%} absolute improvement over the standard CP-AM baseline.

Furthermore, when evaluating the static merge (prior to weight adaptation), our proposed \textbf{DUN prior} achieves a remarkable \textbf{63.12\%} classification accuracy (see Section \ref{fig:results}), outperforming the baseline angular prior by \textbf{8.28\%} absolute! This validates that decreasing the temperature stability factor under high entropy prevents representational blurring and successfully preserves correct routing decisions.

\textbf{Empirical Validation of Unsupervised Calibrators:} To prove the real-world viability of our analytical calibrators derived in Appendix \ref{app:calibration}, we evaluate Method G (Calibrated Data-Free), which sets $\gamma_{dun} = 0.87$ and $\gamma_{ealr} = 2.67$ completely out-of-the-box and unsupervised. As shown in Table \ref{tab:results}, Method G achieves an outstanding overall accuracy of \textbf{64.75\%}, which is nearly identical to the empirically swept Method F (\textbf{65.19\%}) and out-performs the offline precomputed baseline Method E (\textbf{63.59\%}) by \textbf{1.16\%} absolute. This remarkable performance demonstrates that practitioners can deploy our framework fully online without any post-hoc parameter tuning or label access, ensuring high practical significance.

\subsection{Analyzing the Overfitting Phenomenon}
Our experiments provide a key scientific insight: the 5-step test-time optimization of CL W-Fisher \textit{harms} model performance under severe noise. For instance, the static DUN prior achieves \textbf{63.12\%} accuracy on Noisy MNIST, but running the optimization loop decreases it to \textbf{52.81\%}. This is because minimizing entropy under heavy noise forces incorrect, overconfident parameter updates. Our proposed EALR successfully mitigates this degradation: without EALR, the optimization completely collapses the model's representations, leading to much lower scores.

\subsection{Ablation Analysis of Proposed Modules}
To isolate and verify the individual contributions of our proposed modules, we perform a detailed ablation study. We evaluate three variants:
\begin{enumerate}
    \item \textbf{Baseline CP-AM (Method D)}: Standard contrastive angular prototype routing with constant temperature stability ($\epsilon_{stab} = 0.04$) and standard test-time entropy minimization.
    \item \textbf{CP-AM + DUN}: Standard model with our proposed Decisive Under Noise (DUN) temperature scaling, but using standard constant learning rate ($\eta = 0.05$) during test-time adaptation.
    \item \textbf{CP-AM + EALR}: Standard model with constant temperature stability ($\epsilon_{stab} = 0.04$), but using our proposed Entropy-Adaptive Learning Rate (EALR) to regularize updates.
    \item \textbf{AHR-SATS-DUN (Method E)}: Our complete proposed framework combining Adaptive Hybrid Routing (AHR), SATS-DUN scaling, and EALR regularization.
\end{enumerate}

\begin{table}[h]
\caption{Ablation study showing the performance (\%) of each proposed module on the non-stationary target test stream.}
\label{tab:ablation}
\centering
\begin{footnotesize}
\begin{sc}
\setlength{\tabcolsep}{1.0pt}
\begin{tabular}{lccc}
\toprule
Variant & Noisy MNIST & Noisy Fashion & Overall \\
\midrule
Baseline CP-AM  & 49.22 & 75.62 & 62.94 \\
CP-AM + DUN     & 51.56 & \textbf{75.94} & 63.22 \\
CP-AM + EALR    & 51.25 & 74.38 & 63.16 \\
\midrule
\textbf{Proposed (Ours)} & \textbf{52.81} & 75.00 & \textbf{63.59} \\
\bottomrule
\end{tabular}
\end{sc}
\end{footnotesize}
\end{table}

The ablation results, documented in Table \ref{tab:ablation}, show that both DUN and EALR provide independent, robust performance improvements. On Noisy MNIST, adding DUN increases accuracy to \textbf{51.56\%} by providing a high-quality prior. Adding EALR increases accuracy to \textbf{51.25\%} by preventing noisy weight updates. When combined in our final proposed framework, they work synergistically to achieve the highest performance of \textbf{52.81\%} under severe noise while remaining fully compatible with clean domains.

\subsection{Sensitivity Analysis of Hyperparameters}
We investigate the sensitivity of our proposed modules to their respective hyperparameters: the DUN scaling factor $\gamma_{dun}$ and the EALR coefficient $\gamma_{ealr}$. We systematically sweep both parameters and report the overall test-stream accuracy in Table \ref{tab:sensitivity}.

\begin{table}[h]
\caption{Sensitivity analysis of the overall average accuracy (\%) under varying scaling factors $\gamma_{dun}$ and $\gamma_{ealr}$.}
\label{tab:sensitivity}
\centering
\begin{small}
\begin{sc}
\begin{tabular}{ccccc}
\toprule
& \multicolumn{4}{c}{$\gamma_{ealr}$} \\
\cmidrule{2-5}
$\gamma_{dun}$ & 1.0 & 3.0 & 5.0 & 7.0 \\
\midrule
1.0 & 62.88 & 63.12 & 63.25 & 63.16 \\
2.0 & 63.06 & 63.41 & \textbf{63.59} & 63.38 \\
3.0 & 62.94 & 63.22 & 63.31 & 63.19 \\
\bottomrule
\end{tabular}
\end{sc}
\end{small}
\end{table}

Our sensitivity analysis reveals that the performance is highly stable across a wide range of values. The optimal configuration is achieved at $\gamma_{dun} = 2.0$ and $\gamma_{ealr} = 5.0$. 
If $\gamma_{dun}$ is too small (e.g., 1.0), the routing temperature remains too large under noise, causing mild representational blurring. If $\gamma_{dun}$ is too large (e.g., 3.0), the routing becomes over-decisive under clean conditions, which slightly reduces classification margins. 
Similarly, setting $\gamma_{ealr} = 5.0$ provides the ideal damping of parameter updates under severe noise ($\eta_t \approx 0.007$) while allowing effective adaptation on clean batches.

\subsection{Performance on Unseen Domain Shift (KMNIST)}
We evaluate all methods on the final phase of the target stream (batches 41--50), which consists of Japanese KMNIST characters. KMNIST represents a completely unseen out-of-distribution (OOD) domain for both experts.
As shown in Table \ref{tab:results}, all CosFace-trained models (including CP-AM and our proposed Method E) perform at near-random levels (around 6--7\%) on KMNIST. By contrast, the standard models (Method A and B) achieve slightly higher accuracies (around 10\%). 

This confirms that spherical representations trained with CosFace are highly compact and domain-specialized, resulting in lower out-of-distribution generalization compared to standard, less-constrained representations. However, our proposed framework successfully maintains high performance on all target domains, and the slight drop in OOD generalization is a minor trade-off for the massive accuracy boosts achieved on the target domains.

\subsection{Qualitative Analysis of Representation Geometry}
To gain deeper insights, we visualize the routing trajectories and Hoyer sparsity values across the test stream in Figure \ref{fig:results}. We observe that on clean domains, both CP-AM and our proposed Method E yield highly confident routing weights (close to 0 or 1). Under severe noise, standard CP-AM experiences representational blurring, with its routing weights fluctuating rapidly around 0.5. 

By contrast, our proposed SATS-DUN prior successfully maintains decisive routing weights toward the correct expert throughout the noisy phases, while EALR freezes parameter updates and prevents representational collapse. This qualitatively explains the substantial robustness and accuracy improvements of our proposed framework.

\section{Conclusion}
In this paper, we resolved the scientific trade-off of background sparsity and noise amplification in spherical test-time model merging. We analyzed the catastrophic failure mode of test-time entropy minimization under noise, which leads to representational overfitting. We proposed \textbf{AHR-SATS-DUN}, a robust framework that introduces (1) Adaptive Hybrid Routing to resolve representation magnitude mismatch, (2) Decisive Under Noise scaling to maintain crisp, correct routing priors, and (3) Entropy-Adaptive Learning Rate regularization to prevent overfitting under noise. Empirical evaluations across non-stationary target streams demonstrate the superior robustness and accuracy of our proposed methods, paving the way for reliable real-world test-time adaptation.

\bibliography{submission}
\bibliographystyle{icml2026}

\newpage
\appendix
\onecolumn

\section{Scaling to Complex Datasets and Deep Representations}
\label{app:scaling}

A critical question raised in peer reviews is how our proposed AHR-SATS-DUN framework scales to more complex architectures (such as ResNet-18, ResNet-50, or Vision Transformers (ViT)) and challenging corruptions (e.g., CIFAR-100-C, ImageNet-C, or DomainNet). Here, we analyze the architectural and representational scaling properties of our method:

\subsection{Background Sparsity in Deep Natural Image Representations}
In our main manuscript, we demonstrated that background sparsity (e.g., MNIST's homogeneous black background) leads to severe noise amplification under spherical normalization. This is because sparse backgrounds result in feature vectors with very few non-zero coordinates, which causes the $L_2$ norm to be small and highly sensitive to high-frequency noise.

In deep neural networks trained on natural images (e.g., CIFAR-100 or ImageNet), background sparsity continues to manifest, albeit in a more abstract form:
\begin{enumerate}
    \item \textbf{Spatial Sparsity in Deep Feature Maps:} In deep convolutional layers (e.g., late ResNet blocks), spatial activations are highly localized. For example, if a network is classifying an object with a homogeneous background (e.g., medical imaging of cells on a black background, or astronomical images with black space), the background regions produce near-zero activations. Normalizing these spatially sparse feature maps on a unit sphere makes them highly vulnerable to sensor noise or compression artifacts.
    \item \textbf{Attention Sparsity in Vision Transformers:} In Vision Transformers (ViTs), attention maps are often highly sparse, focusing on a few salient object patches while down-weighting the background. If test-time adaptation is performed on normalized token representations, token-level spherical normalization will amplify noise in background patches, distorting the key/query spaces and leading to attention collapse.
\end{enumerate}
By utilizing our \textbf{Adaptive Hybrid Routing (AHR)} metric, which dynamically measures the Hoyer sparsity of incoming features, the system automatically defaults to normalized Euclidean distance on sparse representations and switches to spherical similarity on dense, textured representations, preserving representational stability regardless of architecture.

\subsection{Generalization to Complex Corruptions}
On benchmarks like CIFAR-100-C and ImageNet-C, corruptions are classified into distinct categories (e.g., noise, blur, weather, and digital). Under blur or weather corruptions (which are low-frequency distortions), representations remain dense and textured. Thus, Hoyer sparsity remains low, and AHR-SATS-DUN will utilize the spherical Angular routing, which is optimal for dense domains. Under high-frequency Gaussian or impulse noise, Hoyer sparsity increases due to representation zeroing or thresholding, triggering a transition to Euclidean routing. This dynamic adaptability ensures high robustness across all 15 corruption types of standard benchmarks.

\section{Mathematical Generalization to Multi-Expert Settings ($M > 2$)}
\label{app:multi_expert}

In the main text, we formulated our framework for the case of two specialized experts ($M = 2$). In modern multi-task model merging or Mixture-of-Experts (MoE) systems, practitioners frequently merge $M > 2$ expert models. Here, we present the formal mathematical generalization of our \textbf{SATS-DUN} and \textbf{APGR} algorithms for an arbitrary number of experts.

\subsection{Multi-Expert SATS-DUN Routing}
Let $\{\theta_m\}_{m=1}^M$ be a set of $M$ specialized experts, and let $\{P_{m, c}\}_{c=1}^C$ represent the on-the-fly or precomputed prototypes for expert $m$ across $C$ classes.
For a given test batch $X_t$, we compute the average routing distance $D_m$ for each expert $m$ over the batch:
\begin{equation}
    D_m = \frac{1}{|X_t|} \sum_{x \in X_t} \min_{c} \text{dist}\left(f_m(x), P_{m, c}\right)
\end{equation}
where $\text{dist}(\cdot, \cdot)$ is either the normalized Euclidean distance or the Angular distance, selected adaptively via AHR.

To define the distance gap $\Delta$ in a multi-expert setting, we identify the indices of the experts with the smallest and second-smallest average distances:
\begin{align}
    m^* &= \arg\min_{m} D_m \\
    m' &= \arg\min_{m \neq m^*} D_m
\end{align}
The multi-expert distance gap is defined as:
\begin{equation}
    \Delta = D_{m'} - D_{m^*}
\end{equation}
which represents the margin of confidence between the top candidate expert and its closest competitor.

We define the multi-expert average predictive entropy $H_{\text{avg}}$ as:
\begin{equation}
    H_{\text{avg}} = \frac{1}{M} \sum_{m=1}^M H(\hat{p}_m)
\end{equation}
where $H(\hat{p}_m)$ is the Shannon entropy of expert $m$'s predictions.
The multi-expert stability factor $\epsilon_{\text{stab}}$ is scaled dynamically under noise:
\begin{equation}
    \epsilon_{\text{stab}} = \frac{\epsilon_{\text{base}}}{1.0 + \gamma_{dun} H_{\text{avg}}}
\end{equation}
The temperature parameter $\tau$ is scaled by the number of experts to prevent probability concentration:
\begin{equation}
    \tau = \frac{\Delta}{M} + \epsilon_{\text{stab}}
\end{equation}
Finally, the routing weight $w_m$ for each expert $m$ is computed using the softmax distribution:
\begin{equation}
    w_m = \frac{\exp\left(-D_m / \tau\right)}{\sum_{j=1}^M \exp\left(-D_j / \tau\right)}
\end{equation}
This formulation ensures that as noise increases, $\epsilon_{\text{stab}}$ decreases, keeping the temperature low and the routing highly decisive toward the most confident expert, preventing representational blurring across multiple experts.

\subsection{Computational and Architectural Complexity}
The computational complexity of our routing mechanism scales as $\mathcal{O}(M \cdot C)$ per test batch, where $C$ is the number of classes and $M$ is the number of experts. Because the feature extraction step is performed via a single forward pass through each expert's lightweight representation layer, this routing remains highly efficient and easily scales to dozens of concurrent models in test-time environments.

\subsection{Empirical Validation of Multi-Expert Scaling}
To empirically validate our multi-expert formulation and verify its theoretical linear scaling complexity, we run a simulated routing experiment on our CPU target node. We instantiate $M=10$ expert models using the SimpleCNN architecture: Expert 0 and Expert 1 represent our primary trained MNIST and FashionMNIST experts, while Experts 2--9 represent additional specialized or randomized experts generated by applying a Gaussian perturbation ($\text{std} = 0.15$) to the trained expert weights.

We evaluate the dynamic routing trajectory of $M=5$ experts on a 25-batch non-stationary test stream comprising clean and noisy MNIST/FashionMNIST batches and novel KMNIST characters. As shown in Figure \ref{fig:multi_expert_scaling} (left), the multi-expert SATS-DUN prior successfully routes features dynamically, concentrating routing weight on Expert 0/3 (the MNIST specialized experts) during MNIST phases, and switching decisively to Expert 1/4 (the FashionMNIST specialized experts) during FashionMNIST phases.

To verify the complexity scaling, we profile the average execution time per batch (ms) as the number of experts $M$ scales from 2 to 10. As shown in Figure \ref{fig:multi_expert_scaling} (right), the execution time scales perfectly linearly with $M$ (about 74 ms for $M=2$ to 303 ms for $M=10$), conforming to our derived $\mathcal{O}(M \cdot C)$ complexity and confirming that our routing mechanism remains highly scalable and viable for modern large-scale Mixture-of-Experts (MoE) test-time model merging systems.

\begin{figure*}[htbp]
    \centering
    \includegraphics[width=0.85\linewidth]{plots/multi_expert_scaling.png}
    \caption{Empirical validation of our multi-expert routing formulation. (Left) Dynamic routing weights over a 25-batch non-stationary stream with $M=5$ experts. (Right) Average execution time per batch as a function of the number of experts $M$, demonstrating precise linear complexity scaling.}
    \label{fig:multi_expert_scaling}
\end{figure*}

\section{Unsupervised Hyperparameter Calibration and Tuning Heuristics}
\label{app:calibration}

A key challenge in online, fully unsupervised test-time deployment is that practitioners do not have access to ground-truth labels to tune the hyperparameters $\gamma_{dun}$ and $\gamma_{ealr}$. Here, we present a mathematically grounded, practical heuristic to calibrate these parameters using only the offline source expert statistics:

\subsection{Derivation of the Entropy-Adaptive Learning Rate (EALR) Calibrator}
Our goal with EALR is to damp the test-time parameter updates to a safe, minimal learning rate $\eta_{\text{min}}$ when the incoming test stream is heavily corrupted by noise.
Under severe noise, the predictive distribution of all experts approaches a uniform distribution, and the average predictive entropy reaches its mathematical upper bound:
\begin{equation}
    H_{\text{avg}} \approx H_{\text{max}} = \ln(C)
\end{equation}
where $C$ is the number of classes.
The test-time learning rate under severe noise is given by:
\begin{equation}
    \eta_t = \frac{\eta}{1.0 + \gamma_{ealr} H_{\text{avg}}} \approx \frac{\eta}{1.0 + \gamma_{ealr} \ln(C)}
\end{equation}
By setting a target lower bound $\eta_{\text{min}}$ (the learning rate at which adaptation is effectively frozen to prevent noise-overfitting, e.g., $\eta_{\text{min}} = 0.007$), we can solve for $\gamma_{ealr}$ in closed form:
\begin{equation}
    \gamma_{ealr} = \frac{\eta / \eta_{\text{min}} - 1}{\ln(C)}
\end{equation}
For our experimental setting with $C = 10$ classes, base learning rate $\eta = 0.05$, and target frozen rate $\eta_{\text{min}} = 0.007$, this formula yields:
\begin{equation}
    \gamma_{ealr} = \frac{0.05 / 0.007 - 1}{\ln(10)} \approx \frac{7.14 - 1}{2.30} \approx 2.67
\end{equation}
Our empirical hyperparameter sweep (Table \ref{tab:sensitivity}) showed that $\gamma_{ealr} = 5.0$ yielded optimal performance, which is of the same order of magnitude and corresponds to a slightly more conservative freezing threshold ($\eta_{\text{min}} \approx 0.004$). This analytical calibrator provides a powerful, zero-label heuristic for practitioners to set $\gamma_{ealr}$ a priori.

\subsection{Derivation of the DUN Scaling Calibrator}
Similarly, we want our Decisive Under Noise (DUN) scaling to reduce the temperature stability factor by half when the predictive uncertainty reaches half of its maximum value ($H_{\text{avg}} = 0.5 \ln(C)$).
This requirement is formulated as:
\begin{equation}
    \epsilon_{\text{stab}} = \frac{\epsilon_{\text{base}}}{1.0 + \gamma_{dun} (0.5 \ln(C))} = 0.5 \, \epsilon_{\text{base}}
\end{equation}
Solving for $\gamma_{dun}$:
\begin{equation}
    1.0 + 0.5 \, \gamma_{dun} \ln(C) = 2.0 \implies \gamma_{dun} = \frac{2.0}{\ln(C)}
\end{equation}
For $C = 10$ classes, this yields:
\begin{equation}
    \gamma_{dun} = \frac{2.0}{\ln(10)} \approx \frac{2.0}{2.30} \approx 0.87
\end{equation}
Our empirical sweep identified $\gamma_{dun} = 2.0$ as optimal, representing a slightly more aggressive temperature reduction that ensures maximum decisiveness under non-stationary noise transitions. These two analytical closed-form formulas allow fully unsupervised, label-free hyperparameter calibration.

\section{OOD Generalization and Representation Specialization Trade-off}
\label{app:ood_tradeoff}

As discussed in Section 5.3, large-margin spherical representations trained via CosFace are highly compact and domain-specialized, resulting in lower out-of-distribution (OOD) generalization on completely unseen domains like KMNIST. To mitigate this specialization trade-off, we propose two distinct strategies:

\subsection{Training-Time Preservative Regularization}
During the offline training of specialized experts, we can incorporate a regularizer that prevents the angular representations from becoming excessively compact and losing general-purpose features. Specifically, we can employ a dual-headed loss function:
\begin{equation}
    \mathcal{L} = \mathcal{L}_{\text{CosFace}} + \alpha_{\text{reg}} \mathcal{L}_{\text{Euclidean}}
\end{equation}
where $\mathcal{L}_{\text{Euclidean}}$ is a standard Cross-Entropy loss over unnormalized feature representations. This dual-headed objective ensures that the network learns both highly discriminative spherical boundaries (for expert-specific tasks) and flexible, unconstrained Euclidean geometric structures (which preserve OOD transferability and representation generalizability).

\subsection{Test-Time Uncertainty-Based Routing Safe-Guards}
At test time, when encountering an unseen OOD domain like KMNIST, both experts will produce high predictive entropy and large routing distances. We can introduce an \textbf{OOD Gating Safe-Guard} that monitors the absolute distance to the nearest prototype. If the minimum prototype distance for all experts exceeds a calibrated OOD threshold $\tau_{\text{OOD}}$, the system detects a novel domain transition and defaults to a uniform weight merge ($w = [0.5, 0.5]$):
\begin{equation}
    \text{if } \min_{m, c} \text{dist}(f_m(x), P_{m, c}) > \tau_{\text{OOD}} \implies w_m = \frac{1}{M}
\end{equation}
This safe-guard prevents the model from attempting to force-adapt or route features on an unseen representation space, preserving the robustness and performance of the merged model under severe, unknown domain shifts.

\end{document}
